\section{Upkeep}
\label{sec:reproduced.upkeep}
A refusal that lasted lived on something, and in the record the something was nearly always supplied by the refuser's own people: dues paid in secret, a congregation, a household, a person held in mind.

The plainest record of it is Du Bois's. In the time of slavery, by an account he reprints, there were societies in every Virginia city of any size for caring for the sick and burying the dead. It was unlawful for enslaved people to assemble without a white man present, so the members came to the meeting place by ones and twos, paid the secretary and left. The secretary's book was kept covered up on the bed, and members were known by number. When a member died the others got permits to attend the funeral, and in Richmond they went to the church by ones and twos and sat as near together as they could \autocite[20–21]{dubois1907economic}.

As people were freed the societies came into the open. Philadelphia's Free African Society, founded in April 1787, was formed, its preamble says, “to support one another in sickness, and for the benefit of their widows and fatherless children,” and Du Bois traces the African Methodist church to it \autocite[21]{dubois1907economic}. Its articles set the terms. Each member paid a shilling a month, and after a year's subscription a member in need drew three shillings and ninepence a week. A dead member's widow kept the benefit for as long as she remained his widow. And the society undertook to pay for the education of its dead members' children if they could not attend the free school, and to put them out as apprentices to trades \autocite[21–22]{dubois1907economic}. A society founded to bury its members wrote into its first articles the schooling of the next generation. Upkeep, in this record, is where the formers were paid for.

In Baltimore the societies dated from 1820, and the state later exempted them from its laws forbidding Black people to meet. In 1884 forty of them reported more than 2,100 members. Between them they had buried nearly 1,400 members, paid \$125,000 in sick dues, \$27,000 to widows and more than \$10,700 toward house rent \autocite[92–93]{dubois1907economic}. Those are the numbers the Introduction says a society has always let slide for people: what it costs to keep someone through a sickness, to bury them, to keep their widow housed. The people who paid them wrote them down, first in a secretary's book kept under a bedcover and later in reports read to a public meeting. A burial society refuses nothing. What it does is keep people alive, and together, under an order that would prefer they were neither, and the next section is about what Du Bois saw grow from it.

The congregation did the same work at a larger scale. The church, on Du Bois's account, took over the functions that the other, harshly suppressed social organs had to give up, and became the center of economic life as well as of education, amusement and social intercourse \autocite[24]{dubois1907economic}. It was still doing it more than a century later. The Freedom Schools met in church basements \autocite{payne1995light}, and the Share Croppers' Union held early meetings in churches \autocite{kelley1990hammer}. Nobody in either movement had to build a hall. The congregation already had one, and a membership that turned out every week.

The household carried the rest. A rescue that lasted two to five years, as most of the Oliners' rescues did, was a household's budget, its rooms and its silence for that long, and most rescuers had no one but their own families to share it \autocite{oliner1988altruistic}. The cost of a rescue came due every day the hiding went on, out of the household's ordinary life.

Upkeep is also what keeps the disposition itself from wearing down. Viktor Frankl, who had come through Auschwitz to a camp in Bavaria, marched to a work site in the dark and held to the image of his wife, not knowing whether she was alive \autocite{frankl2006meaning}. Frankl was not refusing an order on that march. His account shows what a refuser needs between decisions, which is an attachment to someone outside the pressure, one that goes on working when the person is absent and might be dead. Nobody assigned it. The camp could not issue it and could not take it.

An exit that pays is the hardest upkeep to supply for oneself, and the record shows how rare it was. In 1956 South Carolina barred public employees from belonging to the NAACP. Septima Clark, a teacher in Charleston, refused to give up her membership and was dismissed. Myles Horton hired her at Highlander to direct its workshops, and the citizenship schools came out of that job \autocite{clark1986ready,horton1990long}. Clark's refusal lasted because a former outside the state had a salary to offer her. Most refusers in the record had nobody with a salary.

The Stasi's officers show the same logic on the operator's side. When the Stasi was dissolved after November 1989, no successor agency took its ninety-one thousand full-time employees; some went into the police and the economy, and a legitimate successor state absorbed the people and paid them \autocite{bpb2023stasi}. The officers had a living on the far side of the order's end, and they formed nothing. What a person lives on decides what they can afford to do next, in both lineages. The refusers in the record mostly had to build their living themselves, out of dues and congregations and households. The Stasi's officers were handed theirs, and taking it was the whole of what was asked of them.

The fifth of the walls is a jobs guarantee for that reason. A wall can pay, for everyone who refuses, the living Highlander paid Clark, without anyone first having to be hired by a movement school. It can't supply the attachment Frankl had or the society that buried its members, and neither can anyone else, because those were the refusers' own. The societies and congregations outlasted every one of their members, and how they did it is the last stage.
